% !TEX root = main.tex
\documentclass[12pt,a4paper]{ltjsarticle}

%% Fonts
\usepackage{lmodern}
\usepackage{luatexja-preset}

%% Layout
\usepackage[top=30mm, bottom=25mm, left=25mm, right=25mm]{geometry}

%% Color / graphics
\usepackage[svgnames,table]{xcolor}
\usepackage{graphicx}

%% Diagrams / charts
\usepackage{tikz}
\usepackage{pgfplots}
\pgfplotsset{compat=1.18}

%% Source code listings
\usepackage{listings}
\renewcommand{\lstlistingname}{ソースコード}

%% Hyperlinks (load near last)
\usepackage{hyperref}

%% Header / footer
\usepackage{fancyhdr}

%% Captions
\usepackage[format=hang, font=small, labelfont=bf]{caption}
\usepackage{subcaption}

%% Mathematics
\usepackage{amsmath}
\usepackage{amssymb}

%% Professional table rules
\usepackage{booktabs}

%% URL
\usepackage{url}

%% Float placement
\usepackage{float}

%% List customization
\usepackage{enumitem}

%% ── Source Code Appearance ───────────────────────────────────
\definecolor{codebg}{RGB}{248,248,248}
\definecolor{codeframe}{RGB}{180,180,180}
\definecolor{kwcolor}{RGB}{0,0,180}
\definecolor{cmcolor}{RGB}{100,100,100}
\definecolor{stcolor}{RGB}{160,0,0}

\lstset{
  backgroundcolor  = \color{codebg},
  basicstyle       = \ttfamily\small,
  keywordstyle     = \color{kwcolor}\bfseries,
  commentstyle     = \color{cmcolor}\itshape,
  stringstyle      = \color{stcolor},
  numbers          = left,
  numberstyle      = \tiny\color{gray},
  numbersep        = 8pt,
  frame            = single,
  rulecolor        = \color{codeframe},
  breaklines       = true,
  breakatwhitespace= false,
  showstringspaces = false,
  tabsize          = 4,
  captionpos       = b,
  xleftmargin      = 2em,
  framexleftmargin = 1.5em,
}

%% ── Page Style ───────────────────────────────────────────────
\pagestyle{fancy}
\fancyhf{}
\rhead{\small プログラミング実験}
\lhead{\small \nouppercase{\leftmark}}
\cfoot{\thepage}
\renewcommand{\headrulewidth}{0.4pt}

\fancypagestyle{titlepage}{
  \fancyhf{}
  \renewcommand{\headrulewidth}{0pt}
}
\fancypagestyle{frontmatter}{
  \fancyhf{}
  \cfoot{\thepage}
  \renewcommand{\headrulewidth}{0pt}
}

%% ── Hyperlink / PDF metadata ─────────────────────────────────
\hypersetup{
  colorlinks = true,
  linkcolor  = NavyBlue,
  citecolor  = NavyBlue,
  urlcolor   = NavyBlue,
  pdftitle   = {レポートタイトル},
  pdfauthor  = {Deka Risman Permana},
  pdfsubject = {プログラミング実験},
  pdfkeywords= {人間情報システム工学科 4年},
}

\newcommand{\HRule}{\rule{\linewidth}{0.5mm}}

\title{}
\date{}
\author{}

%% ============================================================
\begin{document}
%% ============================================================

%% ── Title Page ───────────────────────────────────────────────
\begin{titlepage}
  \thispagestyle{titlepage}
  \centering
  \vspace*{2cm}

  {\large 情報工学実験II}\\[0.5em]
  {\large ハードウェア実験レポート}\\[2em]

  \HRule{}\\[1em]
  {\LARGE \textbf{FPGAを用いたVHDL実験（第１回）}}\\[0.5em]
  \HRule{}

  \vspace{3cm}

  \begin{tabular}{ll}
    \textbf{実験者} & \\[0.8em]
    \textbf{氏　　名} & ：Deka Risman Permana \\[0.4em]
    \textbf{出席番号} & ：4347 \\[1.2em]

    \textbf{共同実験者} & \\[0.8em]
    \textbf{氏　　名} & ：山森　舜水 \\[0.4em]
    \textbf{出席番号} & ：4345 \\[1.2em]

    \textbf{学　　科} & ：人間情報システム工学科 4年 \\[0.4em]
    \textbf{担当教員} & ：縄田　俊則　先生 \\
  \end{tabular}

  \vfill

  \textbf{提出日：}\today

\end{titlepage}

%% ── Table of Contents ────────────────────────────────────────
\newpage
\thispagestyle{frontmatter}
\pagenumbering{roman}
\tableofcontents

\newpage
\pagenumbering{arabic}
\setcounter{page}{1}

%% ── Sections ─────────────────────────────────────────────────

\section{目的}\label{sec:objective}
最初の実験では、Quartusの基本操作を確認しながら、単純な組み合わせ回路を作成してみる。
また、作成した回路をコンポーネントとして別の回路で階層的に利用する方法を習得する。

\section{実験内容}\label{sec:contents}

\subsection{課題 1}\label{subsec:task-1}
最初の課題として、「１つのスイッチの入力をそのまま１つのLEDに出力する」という簡単な回路を作成してみる。
以下がプログラムの例である。

\subsubsection{プログラム}\label{subsubsec:program-1}

\lstinputlisting[language=VHDL, caption=test.vhd]{code/test.vhd}
\newpage
\subsection{課題 2}\label{subsubsec:task-2}
作成した回路は、別の回路でコンポーネント(component、構成要素)として利用することができる。
これは、２入力の論理積の結果を出力する回路である。使用するスイッチやLEDはどれも構わない。動作を確認すること。

\subsubsection{プログラム}\label{subsubsec:program-2}

\lstinputlisting[language=VHDL, caption=andgate.vhd]{code/andgate.vhd}

\lstinputlisting[language=VHDL, caption=andtest.vhd]{code/andtest.vhd}

\subsection{応用課題}
2.1と2.2は動作確認の作業であるが、応用課題として、第２回の課題１の実験をやってみることとした。

\subsubsection{プログラム}\label{subsubsec:program-3}

\lstinputlisting[language=VHDL, caption=adder2bits.vhd]{code/adder2bits.vhd}

\subsection{トップレベルプログラム}\label{subsec:top-level}
全部の課題をコンポーネントとしてまとめて、トップレベルエンティティのプログラムを作成した。

\lstinputlisting[language=VHDL, caption=de2\_top.vhd]{code/de2\_top.vhd}

\subsection{結果}\label{subsec:results}

\subsubsection{課題 1：バッファ回路}
スイッチ SW[0] の入力をそのまま LEDR[0] に出力する回路の動作確認結果を表\ref{tab:result-task1} に示す。

\begin{table}[H]
  \centering
  \caption{課題1 動作確認結果}
  \label{tab:result-task1}
  \begin{tabular}{cc}
    \toprule
    \textbf{SW{[0]}} & \textbf{LEDR{[0]}} \\
    \midrule
    0 & 0 \\
    1 & 1 \\
    \bottomrule
  \end{tabular}
\end{table}

\subsubsection{課題 2：AND ゲート回路}
\texttt{andgate} コンポーネントを2つ用いた回路の動作確認結果を表\ref{tab:result-task2-and1} および表\ref{tab:result-task2-and2} に示す。
LEDR[1] は SW[1] と SW[2] の論理積、LEDR[2] は SW[3] と SW[4] の論理積である。

\begin{table}[H]
  \centering
  \caption{課題2 AND ゲート1（\texttt{and1}）動作確認結果}
  \label{tab:result-task2-and1}
  \begin{tabular}{ccc}
    \toprule
    \textbf{SW{[2]}} & \textbf{SW{[1]}} & \textbf{LEDR{[1]}} \\
    \midrule
    0 & 0 & 0 \\
    0 & 1 & 0 \\
    1 & 0 & 0 \\
    1 & 1 & 1 \\
    \bottomrule
  \end{tabular}
\end{table}

\begin{table}[H]
  \centering
  \caption{課題2 AND ゲート2（\texttt{and2}）動作確認結果}
  \label{tab:result-task2-and2}
  \begin{tabular}{ccc}
    \toprule
    \textbf{SW{[4]}} & \textbf{SW{[3]}} & \textbf{LEDR{[2]}} \\
    \midrule
    0 & 0 & 0 \\
    0 & 1 & 0 \\
    1 & 0 & 0 \\
    1 & 1 & 1 \\
    \bottomrule
  \end{tabular}
\end{table}

\subsubsection{応用課題：2ビット加算回路}
SW[2:1] を入力 $a$、SW[4:3] を入力 $b$ として、それらの和 $c = a + b$ を LEDR[5:3] に出力する回路の動作確認結果を表\ref{tab:result-adder} に示す。
括弧内は10進数表記である。

\begin{table}[H]
  \centering
  \caption{応用課題 2ビット加算回路 動作確認結果}
  \label{tab:result-adder}
  \begin{tabular}{ccc}
    \toprule
    \textbf{$a$：SW{[2:1]}} &
    \textbf{$b$：SW{[4:3]}} &
    \textbf{$c$：LEDR{[5:3]}} \\
    \midrule
    00 (0) & 00 (0) & 000 (0) \\
    00 (0) & 01 (1) & 001 (1) \\
    00 (0) & 10 (2) & 010 (2) \\
    00 (0) & 11 (3) & 011 (3) \\
    01 (1) & 00 (0) & 001 (1) \\
    01 (1) & 01 (1) & 010 (2) \\
    01 (1) & 10 (2) & 011 (3) \\
    01 (1) & 11 (3) & 100 (4) \\
    10 (2) & 00 (0) & 010 (2) \\
    10 (2) & 01 (1) & 011 (3) \\
    10 (2) & 10 (2) & 100 (4) \\
    10 (2) & 11 (3) & 101 (5) \\
    11 (3) & 00 (0) & 011 (3) \\
    11 (3) & 01 (1) & 100 (4) \\
    11 (3) & 10 (2) & 101 (5) \\
    11 (3) & 11 (3) & 110 (6) \\
    \bottomrule
  \end{tabular}
\end{table}

\subsection{ピン設定情報}\label{subsec:pin-settings}

\noindent
\begin{minipage}[t]{0.48\textwidth}
  \centering
  \textbf{スイッチ (SW)}
  \begin{table}[H]
    \centering
    \begin{tabular}{cc}
      \toprule
      \textbf{信号} & \textbf{ピン番号} \\
      \midrule
      SW[7] & PIN\_C13 \\
      SW[6] & PIN\_AC13 \\
      SW[5] & PIN\_AD13 \\
      SW[4] & PIN\_AF14 \\
      SW[3] & PIN\_AE14 \\
      SW[2] & PIN\_P25 \\
      SW[1] & PIN\_N26 \\
      SW[0] & PIN\_N25 \\
      \bottomrule
    \end{tabular}
  \end{table}
\end{minipage}
\hfill
\begin{minipage}[t]{0.48\textwidth}
  \centering
  \textbf{LED (LEDR)}
  \begin{table}[H]
    \centering
    \begin{tabular}{cc}
      \toprule
      \textbf{信号} & \textbf{ピン番号} \\
      \midrule
      LEDR[7]  & PIN\_AC21 \\
      LEDR[6]  & PIN\_AD21 \\
      LEDR[5]  & PIN\_AD23 \\
      LEDR[4]  & PIN\_AD22 \\
      LEDR[3]  & PIN\_AC22 \\
      LEDR[2]  & PIN\_AB21 \\
      LEDR[1]  & PIN\_AF23 \\
      LEDR[0]  & PIN\_AE23 \\
      \bottomrule
    \end{tabular}
  \end{table}
\end{minipage}

\subsection{考察}\label{subsec:discussions}

課題1では，スイッチ SW[0] の入力をそのまま LEDR[0] に出力するバッファ回路を実装した。
ボード上での動作確認の結果，SW[0] をONにすると LEDR[0] が点灯し，OFFにすると消灯することが確認できた。
この結果は表~\ref{tab:result-task1} の真理値表と完全に一致しており，VHDLにおける信号の直接代入（\texttt{led\_out <= sw\_in}）が正しく動作することが確認できた。

課題2では，\texttt{andgate} コンポーネントを2つインスタンス化した \texttt{andtest} 回路を実装した。
ボード上での動作確認の結果，表~\ref{tab:result-task2-and1}・表~\ref{tab:result-task2-and2} に示すとおり，
それぞれのANDゲートは2入力がともにONのときのみLEDが点灯し，それ以外はすべて消灯した。
これはAND演算の定義と一致する。
また，コンポーネントのインスタンス化を用いることで，同一の \texttt{andgate} 回路を2箇所で再利用できることが確認できた

応用課題では，2ビット加算回路 \texttt{adder2bits} を実装した。
ボード上での動作確認の結果，表~\ref{tab:result-adder} に示す全16通りの入力パターンにおいて，
LEDR[5:3] に正しい和が出力されることが確認できた。
入力の最大値は $a = 3,\ b = 3$ であり，その和は $6$（2進数: \texttt{110}）となる。
これは3ビットの出力範囲（0〜7）に収まっており，オーバーフローが発生しないことも確認できた。
\texttt{numeric\_std} ライブラリの \texttt{unsigned} 型を用いることで，
\texttt{std\_logic\_vector} 同士の加算を簡潔に記述できることがわかった。

\newpage
\section{感想}\label{sec:impression}

今回の実験は，FPGAおよびVHDLを用いたハードウェア記述を初めて体験する機会となった。
これまでソフトウェアプログラミングのみを経験してきたため，回路を「記述する」という考え方には最初戸惑いを感じた。
特に，VHDLではすべての信号が並列に動作するという点が，逐次実行を前提とするソフトウェアとの大きな違いであると感じた。
また，Quartusを用いたピン割り当てやコンパイル・ダウンロードの一連の流れを通じて，
FPGAを用いた開発の基本的なワークフローを習得することができた。
コンポーネントを階層的に組み合わせることで，複雑な回路を整理して設計できる点も興味深かった。
今後はより複雑な順序回路にも挑戦してみたいと思う。

\addcontentsline{toc}{section}{参考文献}
\begin{thebibliography}{9}

\bibitem{vhdl-std}
IEEE,
\textit{IEEE Standard VHDL Language Reference Manual},
IEEE Std 1076-2008, IEEE, 2009.
\url{https://ieeexplore.ieee.org/document/4772740}

\bibitem{numeric-std}
IEEE,
\textit{IEEE Standard for VHDL Mathematical Packages (numeric\_std)},
IEEE Std 1076.3-1997, IEEE, 1997.

\end{thebibliography}
\end{document}
